\documentclass[11pt,a4paper]{article}
\usepackage[utf8]{inputenc}
\usepackage[T1]{fontenc}
\usepackage[portuguese]{babel}
\usepackage{geometry}
\geometry{margin=2.5cm}
\usepackage{listings}
\usepackage{xcolor}
\usepackage{amsmath}
\usepackage{amssymb}
\usepackage{hyperref}
\usepackage{enumitem}
\usepackage{tcolorbox}
\usepackage{tabularx}
\usepackage{booktabs}

\definecolor{codebg}{rgb}{0.95,0.95,0.95}
\definecolor{codegreen}{rgb}{0.1,0.5,0.1}
\definecolor{codepurple}{rgb}{0.5,0.1,0.5}
\definecolor{codegray}{rgb}{0.5,0.5,0.5}

\lstdefinestyle{python}{
    backgroundcolor=\color{codebg},
    commentstyle=\color{codegray}\itshape,
    keywordstyle=\color{codepurple}\bfseries,
    stringstyle=\color{codegreen},
    basicstyle=\ttfamily\small,
    breaklines=true,
    frame=single,
    language=Python,
    showstringspaces=false,
    tabsize=4,
    literate={á}{{\'a}}1 {ã}{{\~a}}1 {é}{{\'e}}1 {í}{{\'i}}1 {ó}{{\'o}}1 {õ}{{\~o}}1 {ú}{{\'u}}1 {ç}{{\c{c}}}1 {Á}{{\'A}}1 {Ã}{{\~A}}1 {É}{{\'E}}1 {Í}{{\'I}}1 {Ó}{{\'O}}1 {Õ}{{\~O}}1 {Ú}{{\'U}}1 {Ç}{{\c{C}}}1 {â}{{\^a}}1 {ê}{{\^e}}1 {ô}{{\^o}}1 {û}{{\^u}}1 {Â}{{\^A}}1 {Ê}{{\^E}}1 {Ô}{{\^O}}1 {Û}{{\^U}}1 {à}{{\`a}}1 {è}{{\`e}}1 {ì}{{\`i}}1 {ò}{{\`o}}1 {ù}{{\`u}}1
}
\lstset{style=python}

\newtcolorbox{solucao}{
  colback=yellow!5!white,
  colframe=orange!60!black,
  title=Solução proposta,
  fonttitle=\bfseries
}

\title{\textbf{Data Analysis with Python 2024--2025}\\Parte 3: Matplotlib\\\large Soluções dos Exercícios}
\author{Soluções independentes elaboradas para fins de estudo, com base nos enunciados originais}
\date{}

\begin{document}
\maketitle

\section*{Nota Importante}
Este material é uma adaptação das soluções independentes para os exercícios baseados no curso \textit{Data Analysis with Python 2024-2025}, oferecido pela \textbf{The University of Helsinki MOOC Center}. As soluções abaixo não são o gabarito oficial do curso.

\tableofcontents
\newpage

\section{Soluções - Capítulo 3 (Matplotlib)}

\subsection*{Exercício 3.9 -- Vários gráficos}
\begin{solucao}
\begin{lstlisting}
import matplotlib.pyplot as plt

def main():
    x1 = [2, 4, 6, 7]
    y1 = [4, 3, 5, 1]
    
    x2 = [1, 2, 3, 4]
    y2 = [4, 2, 3, 1]
    
    plt.plot(x1, y1)
    plt.plot(x2, y2)
    
    plt.title("Dois graficos no mesmo conjunto de eixos")
    plt.xlabel("Eixo x")
    plt.ylabel("Eixo y")
    
    plt.show()

if __name__ == "__main__":
    main()
\end{lstlisting}
\end{solucao}

\subsection*{Exercício 3.10 -- Subfiguras}
\begin{solucao}
\begin{lstlisting}
import numpy as np
import matplotlib.pyplot as plt

def subfigures(a):
    fig, ax = plt.subplots(1, 2)
    
    ax[0].plot(a[:, 0], a[:, 1])
    
    ax[1].scatter(
        a[:, 0],
        a[:, 1],
        c=a[:, 2],
        s=a[:, 3]
    )
    
    plt.show()

def main():
    a = np.array([
        [1, 2, 10, 20],
        [2, 3, 20, 40],
        [3, 1, 30, 60],
        [4, 5, 40, 80]
    ])
    subfigures(a)

if __name__ == "__main__":
    main()
\end{lstlisting}
\end{solucao}

\vspace{2em}
\noindent\textbf{Créditos:} Este conteúdo foi originalmente criado pela \textit{The University of Helsinki MOOC Center} para o curso \textit{Data Analysis with Python}.
\end{document}
